\documentclass[10pt,a4paper]{amsart}
\usepackage{amsmath,amssymb,amsthm,graphicx}
% Retained minimal preamble for the Stacks Project chapter sites-cohomology.tex (Cohomology on Sites).
% Verbatim excerpt of the chapter preamble: Xy-pic setup (native preamble lines 14-18) and the native
% theorem declarations with the mathematical macros (native preamble lines 154-217). Copied mechanically;
% nothing added, removed or redefined. The wrapper class and the amsmath/amssymb/amsthm/graphicx base
% packages are supplied by the assembler, so the class/externaldocument part of the native preamble is omitted.
% For commutative diagrams we use Xy-pic
\usepackage[all]{xy}

% We use 2cell for 2-commutative diagrams.
\xyoption{2cell}
% Theorem environments.
%
\theoremstyle{plain}
\newtheorem{theorem}[subsection]{Theorem}
\newtheorem{proposition}[subsection]{Proposition}
\newtheorem{lemma}[subsection]{Lemma}

\theoremstyle{definition}
\newtheorem{definition}[subsection]{Definition}
\newtheorem{example}[subsection]{Example}
\newtheorem{exercise}[subsection]{Exercise}
\newtheorem{situation}[subsection]{Situation}

\theoremstyle{remark}
\newtheorem{remark}[subsection]{Remark}
\newtheorem{remarks}[subsection]{Remarks}

\numberwithin{equation}{subsection}

% Macros
%
\def\lim{\mathop{\mathrm{lim}}\nolimits}
\def\colim{\mathop{\mathrm{colim}}\nolimits}
\def\Spec{\mathop{\mathrm{Spec}}}
\def\Hom{\mathop{\mathrm{Hom}}\nolimits}
\def\Ext{\mathop{\mathrm{Ext}}\nolimits}
\def\SheafHom{\mathop{\mathcal{H}\!\mathit{om}}\nolimits}
\def\SheafExt{\mathop{\mathcal{E}\!\mathit{xt}}\nolimits}
\def\Sch{\mathit{Sch}}
\def\Mor{\mathop{\mathrm{Mor}}\nolimits}
\def\Ob{\mathop{\mathrm{Ob}}\nolimits}
\def\Sh{\mathop{\mathit{Sh}}\nolimits}
\def\NL{\mathop{N\!L}\nolimits}
\def\CH{\mathop{\mathrm{CH}}\nolimits}
\def\proetale{{pro\text{-}\acute{e}tale}}
\def\etale{{\acute{e}tale}}
\def\QCoh{\mathit{QCoh}}
\def\Ker{\mathop{\mathrm{Ker}}}
\def\Im{\mathop{\mathrm{Im}}}
\def\Coker{\mathop{\mathrm{Coker}}}
\def\Coim{\mathop{\mathrm{Coim}}}

% Boxtimes
%
\DeclareMathSymbol{\boxtimes}{\mathbin}{AMSa}{"02}

%
% Macros for moduli stacks/spaces
%
\def\QCohstack{\mathcal{QC}\!\mathit{oh}}
\def\Cohstack{\mathcal{C}\!\mathit{oh}}
\def\Spacesstack{\mathcal{S}\!\mathit{paces}}
\def\Quotfunctor{\mathrm{Quot}}
\def\Hilbfunctor{\mathrm{Hilb}}
\def\Curvesstack{\mathcal{C}\!\mathit{urves}}
\def\Polarizedstack{\mathcal{P}\!\mathit{olarized}}
\def\Complexesstack{\mathcal{C}\!\mathit{omplexes}}
% \Pic is the operator that assigns to X its picard group, usage \Pic(X)
% \Picardstack_{X/B} denotes the Picard stack of X over B
% \Picardfunctor_{X/B} denotes the Picard functor of X over B
\def\Pic{\mathop{\mathrm{Pic}}\nolimits}
\def\Picardstack{\mathcal{P}\!\mathit{ic}}
\def\Picardfunctor{\mathrm{Pic}}
\def\Deformationcategory{\mathcal{D}\!\mathit{ef}}

% --- appended 2026-10-12 (Chief): the project own definitions of the reference/slogan/history
% display environments, copied mechanically from the shared Stacks preamble lines 8-12
% (\usepackage{verbatim} plus the three \newenvironment lines, which the chapter preamble keeps
% verbatim from the project preamble). Nothing added, removed or redefined: these environments
% carry no mathematical content and the project itself typesets them as comments. Without these
% three lines a counted statement whose own text uses them fails with Environment slogan undefined.
\usepackage{verbatim}
\newenvironment{reference}{\comment}{\endcomment}
\newenvironment{slogan}{\comment}{\endcomment}
\newenvironment{history}{\comment}{\endcomment}
\usepackage{cleveref}
\crefname{originallemma}{original lemma}{original lemmas}
\Crefname{originallemma}{Original lemma}{Original lemmas}
\makeatletter
\newlabel{derived-lemma-right-derived-exact-functor}{{015F}{source}{}{sourcepointer-1}{}}
\newlabel{derived-lemma-right-derived-exact-functor@cref}{{[originallemma][0][]015F}{[1][0][]source}{}{}{}}
\makeatother
\crefname{originalsection}{original section}{original sections}
\Crefname{originalsection}{Original section}{Original sections}
\makeatletter
\newlabel{homology-section-functors}{{010M}{source}{}{sourcepointer-2}{}}
\newlabel{homology-section-functors@cref}{{[originalsection][0][]010M}{[1][0][]source}{}{}{}}
\makeatother
\title{Sections over a weakly contractible object are exact, acyclic and hit torsors}
\author{The Stacks Project Authors}
\date{stacks-project:sites-cohomology.tex:local-snapshot-20261011; GNU FDL 1.2 or later}
\begin{document}
\maketitle
\noindent\textit{Editorial note.} All delivered mathematical text is the author's own, copied from the pinned Stacks Project chapter file sites-cohomology.tex; no statement or proof step was rewritten, reordered or removed.\par\medskip
\noindent\textit{Editorial note.} The counted proof is the author's complete argument: the proof is written out in full: the two statements read off the definition and the vanishing of the higher derived functors of the exact sections functor.\par\medskip
\noindent\textit{Editorial note.} The label lemma-w-contractible names the counted statement, and the proof environment immediately following it in the source is delivered verbatim as its proof.\par\medskip
\noindent\textit{Editorial note.} The counted statement is carried by the author's own lemma environment in the chapter source.\par\medskip
\noindent\textit{Editorial note.} The statement and proof are copied from the pinned chapter file; the two cited inputs are kept as external pointers with their canonical Stacks tags.\par\medskip
\noindent\textit{Allowed standard prerequisites (editorial).}\begin{itemize}
\item Weakly contractible objects of a site. Same chapter, the definition of a weakly contractible object through the exactness of the sections functor and the triviality of the torsors over it. Supplies the object whose three properties are asserted.
\item Right derived functors of an exact functor vanish. Derived Categories, the lemma stating that the right derived functors of an exact functor vanish in positive degrees. Supplies the vanishing of the higher cohomology groups in part (2).
\end{itemize}\medskip
\ifdefined\hypertarget\hypertarget{sourcepointer-1}{}\fi
\noindent\textit{Source reference.} Original lemma 015F: sources/stacks-chow/tags.txt, source lines 1506--1506. Prerequisite pointer invoked at proof line(s) 13790; it names a standard fact of a sibling chapter that the counted argument uses, and every step of the counted proof built on it is retained verbatim in the delivered body.\par\medskip
\ifdefined\hypertarget\hypertarget{sourcepointer-2}{}\fi
\noindent\textit{Source reference.} Original section 010M: sources/stacks-chow/tags.txt, source lines 1332--1332. Prerequisite pointer invoked at proof line(s) 13788; it names a standard fact of a sibling chapter that the counted argument uses, and every step of the counted proof built on it is retained verbatim in the delivered body.\par\medskip
\begin{lemma}
\label{lemma-w-contractible}
Let $\mathcal{C}$ be a site. Let $U$ be a weakly contractible
object of $\mathcal{C}$. Then
\begin{enumerate}
\item the functor $\mathcal{F} \mapsto \mathcal{F}(U)$ is an exact
functor $\textit{Ab}(\mathcal{C}) \to \textit{Ab}$,
\item $H^p(U, \mathcal{F}) = 0$
for every abelian sheaf $\mathcal{F}$ and all $p \geq 1$, and
\item for any sheaf of groups $\mathcal{G}$ any $\mathcal{G}$-torsor
has a section over $U$.
\end{enumerate}
\end{lemma}

\begin{proof}
The first statement follows immediately from the definition
(see also Homology, Section \ref{homology-section-functors}).
The higher derived functors vanish by
Derived Categories, Lemma \ref{derived-lemma-right-derived-exact-functor}.
Let $\mathcal{F}$ be a $\mathcal{G}$-torsor. Then $\mathcal{F} \to *$
is a surjective map of sheaves. Hence (3) follows from the
definition as well.
\end{proof}

\end{document}
